% 31-05(31쪽) — 축이 직선 x=3 인 이차함수 y=f(x) 의 그래프와 두 점 A(3, 0), B(6, f(6)) 을 지나는 직선 AB. 스캔 화질이 낮아 TikZ 로 다시 그림.
% 원문에 있는 것만 옮긴다: 포물선 · 직선 AB · B 에서 x축으로 내린 점선 · 라벨 O, x, y, A, B, 3, 6, y=f(x). 세로 눈금값은 원문에 없다.
\begin{tikzpicture}[x=0.50cm, y=0.30cm, line width=0.6pt]
  \draw[->, >=stealth, line width=0.7pt] (-0.9,0) -- (7.4,0) node[below=1pt, font=\small] {$x$};
  \draw[->, >=stealth, line width=0.7pt] (0,-1.7) -- (0,6.1) node[left=1pt, font=\small] {$y$};
  \draw[densely dashed, line width=0.4pt] (6,0) -- (6,4);
  \draw[line width=0.6pt, domain=0.0:6.55, samples=90, smooth] plot (\x, {4*(\x-3)*(\x-3)/9});
  \draw[line width=0.6pt] (2.05,-1.27) -- (6.6,4.8);
  \node[below left=-2pt and -2pt, font=\small] at (0,0) {$\pt{O}$};
  \node[above left=-1pt and -2pt, font=\small] at (3,0) {$\pt{A}$};
  \node[above left=-2pt and -1pt, font=\small] at (6,4) {$\pt{B}$};
  \node[below=1pt, font=\small] at (3,0) {$3$};
  \node[below=1pt, font=\small] at (6,0) {$6$};
  \node[font=\small] at (3.5,5.3) {$y=f(x)$};
\end{tikzpicture}
